\section{Conclusion}
\label{sec:conclusion}

Reliable, real-time traffic anomaly detection in intelligent transportation systems is frequently hindered by decentralized data silos, extreme class imbalance, and the strict hardware constraints of vehicular edge nodes. This paper presented FedRouteAI, a federated learning framework that pairs a deployment-aware inference pipeline with an adaptive aggregation strategy to address these challenges. The core methodological contributions include a federated architecture that preserves data privacy across distinct regional clients, a historical feature engineering module for localized normalization, the four-factor AdaptiveFedAvg aggregation algorithm, and a deployment-aware inference mechanism relying on zero-imputation and dynamic Out-of-Distribution (OOD) warnings.

Under the evaluated synthetic dataset, the federated approaches achieved higher generalization performance than the centralized baseline. AdaptiveFedAvg's performance was comparable to standard FedAvg, with only marginal differences in global accuracy and Macro F1, which suggests that standard, dataset-proportional aggregation already performs robustly in stable deployment environments. Systems-level benchmarking additionally shows that the lightweight Multi-Layer Perceptron architecture operates with low inference latency and minimal communication overhead, meeting the latency and communication requirements evaluated in this study and remaining compatible with real-time edge deployment scenarios.

Despite these promising indications, several limitations must be acknowledged. The experimental evaluation relied entirely on a synthetic dataset and a simulated federated environment, which limits the validation of the system under true physical conditions, and it was constrained to a single MLP architecture that was not tested against the severe, erratic packet drops typical of physical edge sensors. Future work should evaluate the framework using larger and more heterogeneous federated deployments, investigate more advanced learning architectures (e.g., Graph Neural Networks or Transformer-based models), and validate the system under real roadside infrastructure.
